% R592: whole EGA I Chapter0 3.3.2–3; source-faithful translation.
% OpenAI Codex, GPT-6.1 Sol, Ultra; no human review claimed.
\begin{env}[3.3.2]
\label{0.3.3.2-fr}
Con las mismas notaciones, sea $\mathcal{G}_\lambda$ un segundo haz sobre
$U_\lambda$ (para todo $\lambda\in L$), con valores en $\boldsymbol{K}$, y supongamos
dado para todo par $(\lambda,\mu)$ un isomorfismo
\[
  \omega_{\lambda\mu}:
  \mathcal{G}_\mu|(U_\lambda\cap U_\mu)
  \xrightarrow{\sim}
  \mathcal{G}_\lambda|(U_\lambda\cap U_\mu),
\]
estos isomorfismos verifican la condición de pegado. Supongamos finalmente dado
para cada $\lambda$ un morfismo
$u_\lambda:\mathcal{F}_\lambda\to\mathcal{G}_\lambda$, y que los diagramas
\begin{equation}
\begin{tikzcd}[column sep=huge,row sep=large]
\mathcal{F}_\mu|(U_\lambda\cap U_\mu)
  \arrow[r,"u_\mu"] \arrow[d]
  & \mathcal{G}_\mu|(U_\lambda\cap U_\mu) \arrow[d] \\
\mathcal{F}_\lambda|(U_\lambda\cap U_\mu)
  \arrow[r,"u_\lambda"']
  & \mathcal{G}_\lambda|(U_\lambda\cap U_\mu)
\end{tikzcd}
\tag{3.3.2.1}\label{0.3.3.2.1-fr}
\end{equation}
sean conmutativos. Entonces, si $\mathcal{G}$ se obtiene por pegado de los
$\mathcal{G}_\lambda$ mediante los $\omega_{\lambda\mu}$, existe un
morfismo $u:\mathcal{F}\to\mathcal{G}$ y uno solo tal que los diagramas

\oldpage[0\textsubscript{I}]{29}

\[
\begin{tikzcd}[column sep=huge,row sep=large]
\mathcal{F}|U_\lambda
  \arrow[r,"u|U_\lambda"] \arrow[d]
  & \mathcal{G}|U_\lambda \arrow[d] \\
\mathcal{F}_\lambda
  \arrow[r,"u_\lambda"']
  & \mathcal{G}_\lambda
\end{tikzcd}
\]
sean conmutativos~; esto se deduce inmediatamente de (3.2.3). La correspondencia entre
la familia $(u_\lambda)$ y $u$ es una biyección funtorial del subconjunto de
$\prod_\lambda\operatorname{Hom}(\mathcal{F}_\lambda,\mathcal{G}_\lambda)$
que verifica las condiciones (3.3.2.1) sobre
$\operatorname{Hom}(\mathcal{F},\mathcal{G})$.
\end{env}

\begin{env}[3.3.3]
\label{0.3.3.3-fr}
Con las notaciones de (3.3.1), sea $V$ un abierto de $X$~; es inmediato que
las restricciones a $V\cap U_\lambda\cap U_\mu$ de los $\theta_{\lambda\mu}$
satisfacen la condición de pegado para los haces inducidos
$\mathcal{F}_\lambda|(V\cap U_\lambda)$ y que el haz sobre $V$ obtenido por
pegado de estos últimos se identifica canónicamente con $\mathcal{F}|V$.
\end{env}
